\acresetall
% Introduction (outline brief 6.1; mirrors RQ1 L35-102, seven paragraphs to seven). The abstract expands DCAT,
% DCAT-AP and LLM as RQ1's does; \acresetall expands them again here (verify, not printed: RQ1 main.tex convention,
% readme-notes question 1). Describes the tool at 5ffd56e (DCAT-state pinning, rebase). No development-run number.
\section{Introduction}\label{sec:intro}

The W3C \ac{DCAT} is a vocabulary for interoperability between data catalogs~\cite{w3cdcat3}, and the
\ac{DCAT-AP} adds obligations, cardinalities and controlled vocabularies~\cite{semicdcatap}. Domain extensions
build on \acs{DCAT-AP} and inherit its classes, properties and obligations: GeoDCAT-AP and StatDCAT-AP for
geospatial and statistical data~\cite{geodcatap10,statdcatap}, mobilityDCAT-AP for mobility
data~\cite{scrocca2025mobilitydcat}, HealthDCAT-AP for the \ac{EHDS}~\cite{derycke2025healthdcat} and DCAT-AP-JRC
for research data~\cite{perego2016jrc}. What an extension adds is a set of changes to its base: existing
properties used where the base states nothing, base constraints tightened and terms minted in the extension's own
namespace~\cite{semicreuse2026}.
% src for the three kinds of change: @5ffd56e docs/conformance.md:12-14, :45 (RG lines 376-379, 471-474)

Defining such an extension is a requirements problem before it is a modeling problem~\cite{anon2027rq1}; once
the requirements are settled, it is a modeling problem under the constraints the base imposes. A requirement is
met by the base as it stands, by a tightened base constraint, by a reused term or by a new one; mobilityDCAT-AP,
for example, makes the update frequency mandatory where its base has it optional~\cite{mobilitydcatspec}.
% src: RQ1 L59-61 (the companion's own example)
\todo{author: keep the companion's mobilityDCAT-AP example (a test case) or use the generic wording ``raising an
optional property to mandatory'' (Q17)}
The guidance of the \ac{SEMIC} orders these choices, rules out loosening a base constraint, ``as the result would
no longer be compliant with DCAT-AP'', and expects use cases to justify new elements and tightened
constraints~\cite{semicreuse2026}. \verify{verbatim against the guidelines}
% src: @5ffd56e docs/conformance.md:12-14 (RG lines 376-379) and :47 (RG lines 262-264)
The base itself profiles a DCAT version as it stood when the release was made: \acs{DCAT-AP}~1.1 builds on % src: @5ffd56e B/data/base-profiles/dcat-ap-1.1/profile.yaml:33
DCAT~2014 and \acs{DCAT-AP}~3.0.1 on DCAT~3. \verify{release quotes} % src: @5ffd56e B/data/base-profiles/dcat-ap-3.0.1/profile.yaml:43

Existing authoring tools and generators help with the modeling, not with this derivation
(Sect.~\ref{sec:related}). Dataspecer generates \ac{DSV} and \ac{SHACL} from profiles of profiles that people
author \verify{new bib entry: Dataspecer}, but in a reading of its code it uses no language model and does not
check that a profile refines its base; generators based on \acp{LLM} write ontologies from competency questions or
relational schemas~\cite{lippolis2025ontology,fathallah2026rigor}, and shape generators are given the
terms \verify{new bib entry: NL2SHACL-Bench}. Of those we examined, none derives an extension of a pinned base
release from requirements and records, for every element, the requirement it serves and the base value it
changes. \verify{each ``none'' claim}
% src: docs/research/B1-dataspecer.md:29 (finding 7: no refinement check); visual-profile-editor
% docs/dataspecer-spike.md:9-30, 67-75 (SHACL from profiles; no LLM features, medium confidence: code reading only)

Such records are needed after the first release too, because a base does not stand still. \acs{DCAT-AP}~3.0.1
relaxed three cardinalities of 3.0.0 \verify{new bib entry: DCAT-AP 3.0.1}, and 68 issues on \acs{DCAT-AP}, % src: @5ffd56e docs/conformance.md:22-23 (three cardinalities, from the 3.0.1 CHANGELOG); RQ1 L71-73 (68 issues)
StatDCAT-AP and GeoDCAT-AP were reviewed for maintenance in less than nine months~\cite{dragan2019maintenance}. % src: RQ1 L71-73
\verify{3.0.1 CHANGELOG}
When the base moves to another release, and possibly to another DCAT version, an extension has to find the elements
that rest on what changed, and a profile that names neither the requirement behind an element nor the base value
it tightened cannot show it. With such a record, the rebase of Sect.~\ref{sec:impl} compares each change under
the old and the new base.
% src: @5ffd56e B/engine/rebase.py:73-180; B/engine/pipeline.py:557-609; docs/dcat-versions.md:185-210

This paper addresses the gap with a pipeline in which an \acs{LLM} acts as a bounded decider: for each
requirement atom it decides whether the base already meets it and otherwise fills typed change records from
retrieved candidate terms, while program code mints identifiers and compiles the extension. Four pillars shape % src: @5ffd56e docs/design.md:5-10 (the design goals; outline section 2)
the design: (1)~reuse first and minimal, (2)~traced, with an obligation the atom's wording licenses,
(3)~decided by the model, with the final say for a person and every step open to inspection, and
(4)~conformant to the base; Sect.~\ref{sec:problem} turns them into the criteria E1--E4. The companion workflow
ends with needs that cite their passages; this pipeline starts from atoms and ends with changes that cite them,
each with a verbatim quote of its atoms, the base value beside each new one and the decisions behind it, so that
when an atom or the base changes a reviewer can find the changes that rest on it. Program code checks every
change and the compiled extension, a model only judges whether the atoms support a change, and only a person
releases the extension; automatic acceptance is a decision policy, not a claim that the checks establish
correctness.
% src: @5ffd56e B/engine/prompts.py:24-72; B/engine/verify.py:482-483; B/engine/pipeline.py:221-267, 611-640;
% docs/design.md:14-31

We ask, for the modeling part of profile engineering: \emph{How much of what a published DCAT-AP extension adds
to its base does a pipeline reconstruct from the requirement atoms behind it that reuses before it mints and
keeps every change reviewable from the atoms it covers, its relation to the base and its decisions, while the
extension stays conformant to the base, and what do retrieval, repair and the judge contribute?} We compare the
full pipeline with a rules baseline and a closed-book model, and remove retrieval, repair and the judge in turn.
\todo{conditions: final set (Q7)}
% src: @5ffd56e B/contracts.py:592-596 (the switches retrieval, repair_rounds, attribution_check)

This paper makes three contributions:
\begin{itemize}
\item a formulation of an extension as a justified delta over a pinned \acs{DCAT-AP} release and the DCAT version
  it builds on, with six change kinds, a justification contract and a legal-refinement order, and a pipeline that % src: @5ffd56e B/contracts.py:39-46 (ChangeKind)
  implements it, grounded step by step in requirements-engineering and ontology-generation work
  (Sects.~\ref{sec:problem} and~\ref{sec:pipeline});
\item reference deltas for the five published extensions of the companion benchmark, defined as each published % src: RQ1 L381-388
  profile minus its base and matched to the generated changes element by element (Sect.~\ref{sec:benchmark});
  \todo{reference deltas: build, align and freeze before this claim stands}
\item an evaluation design in two views, how much of the reference delta the output recovers and how much of the
  output the reference accounts for, with reuse, conformance, review load and cost (Sects.~\ref{sec:setup}
  and~\ref{sec:results}). \todo{switch ``an evaluation design'' to ``an evaluation'' when results exist}
\end{itemize}
